\begin{center}
  \large\textbf{ABSTRACT}
\end{center}

\addcontentsline{toc}{chapter}{ABSTRACT}

\vspace{2ex}

\begingroup
% Menghilangkan padding
\setlength{\tabcolsep}{0pt}

\noindent
\begin{tabularx}{\textwidth}{l >{\centering}m{3em} X}
  \emph{Name}     & : & \name{}         \\

  \emph{Title}    & : & \engtatitle{}   \\

  \emph{Advisors} & : & 1. \advisor{}   \\
                  &   & 2. \coadvisor{} \\
\end{tabularx}
\endgroup

% Ubah paragraf berikut dengan abstrak dari tugas akhir dalam Bahasa Inggris
\emph{
  This study designs and evaluates a \textit{Credit Risk Scoring} system based on \textit{Machine Learning Operations} (MLOps) to support adaptive, measurable, and operationally sustainable credit risk assessment. A synthetic dataset of 10,000 samples is used to represent retail borrower profiles, then processed through \textit{Weight of Evidence} (WoE) transformation, \textit{Information Value} (IV)-based feature selection, and class imbalance handling. The main predictive model is built using XGBoost, while the MLOps pipeline integrates experiment tracking, model registry, API deployment, and automated \textit{Continuous Training} executed via schedule and event-based (new data accumulation) triggers. The evaluation results indicate that this approach produces an accurate model, improves interpretability through SHAP, and is supported by an automation workflow that consistently ensures model freshness, thereby enhancing reliability and implementation readiness in financial institutions.}

% Ubah kata-kata berikut dengan kata kunci dari tugas akhir dalam Bahasa Inggris
\emph{Keywords}: \emph{Credit risk scoring}, MLOps, XGBoost, WoE, SHAP.
